\documentclass[10pt,a4paper]{amsart}
\usepackage[margin=19mm]{geometry}
\usepackage{amsmath,amssymb,mathtools}
\usepackage[scr=rsfs,cal=euler]{mathalfa}
\usepackage{xfrac}
\usepackage[unicode,hidelinks,bookmarks=false]{hyperref}
\newcommand{\ie}{i.e.\ }
\newcommand{\cf}{cf.\ }
\newcommand{\resp}{respectively\ }
\newcommand{\Subsection}{\subsection}
\providecommand{\nobreakdash}{\nobreak\hbox{-}}
\setlength{\emergencystretch}{2em}
\multlinegap0pt
\usepackage{tikz}
\usepackage{xcolor}
\usepackage{amssymb}
\usepackage[shortlabels]{enumitem}
\newtheorem{lemma}{Lemma}
\title{An elementary approach to Sun Kim's general theta function identities}
\author{Dazhao Tang}
\date{Source: arXiv:2609.09751v1 (CC BY 4.0)}
\begin{document}
\maketitle

\noindent\emph{Source and licence.} An elementary approach to Sun Kim's general theta function identities. Version arXiv:2609.09751v1 (CC BY 4.0), distributed by arXiv under the Creative Commons Attribution 4.0 International licence, http://creativecommons.org/licenses/by/4.0/, as recorded in the arXiv licence metadata for this exact version and shown on the abstract page https://arxiv.org/abs/2609.09751v1. Full licence notice: \texttt{licenses/arxiv-2609.09751-CC-BY-4.0.txt}.\par\smallskip

\noindent\textit{Editorial scope.} Counted unit: the article's lemma Fund-Lem-1 (source0.tex lines 278--297). The article's complete original proof of it is delivered with it (source0.tex lines 299--389, one \texttt{proof} environment of 91 lines). No mathematics is written, repaired or re-derived: the statement and the proof are the article's own lines, in the article's own notation, and no retained line is substituted or re-flowed.  Retained uncounted as the source-local context the proof needs: lines 159--162.  Nothing was omitted from any retained span, so the package carries no omission record.  No retained line contains a citation, so the wrapper carries no bibliography and no bibliography entry.  No retained line contains a figure environment, an image-inclusion command or drawing code, so no image is delivered and the package declares no asset dependency.  Editorial formatting only, in addition: an AMS \texttt{amsart} preamble that restates the article's own declarations and macros used by the retained lines, namely the environments \texttt{lemma} on separate counters, together with the standard packages those lines need.  The retained environments print in this document as Lemma 1 in order of appearance, whereas the article prints the same items with the numbering of its own full document.  The complete native source of the article as distributed in the arXiv e-print for this exact version is bundled unchanged as \texttt{sources/209850/source0.tex}.
\begin{align}
\theta(-c,-ac,-bc,-abc;q) &-\theta(c,ac,bc,abc;q)\nonumber\\
 &\qquad\qquad=c\!\;\theta\big({-}1,-a,-b,-abc^2;q\big),\label{War-gene}
\end{align}

\begin{lemma}\label{Fund-Lem-1}
Let
\begin{align}
H(X,Y)=\theta\big(qX;q^2\big)\theta\big({-}qY;q^2\big)+\theta\big({-}qX;q^2)
\theta\big(qY;q^2\big).\label{Def-H}
\end{align}
Then
\begin{align}
\mathscr{D}\big(x_1^2,x_2^2,x_3^2\big) &-\mathscr{S}\big(x_1^2,x_2^2,x_3^2\big)
\mathscr{D}\big(qx_1^2,qx_2^2,qx_3^2\big)\nonumber\\
 &\qquad\qquad=-\dfrac{c^2z^2}{q^2}\dfrac{\theta\big(a^2,b^2;q^2\big)}{\theta(a,b;q^2)}
C(r,z),\label{Fund-iden-1}
\end{align}
where $r=d/c$, $z=abc^2d^2$, and
\begin{align}
C(r,z) &=r^2\theta\big({-}z/r^2;q^2\big)H\big(zr^2,z\big)+\theta\big({-}zr^2;q^2\big)
H\big(z,z/r^2\big)\nonumber\\
 &\quad+r\!\:\theta\big({-}z;q^2\big)H\big(z/r^2,zr^2\big).\label{Def-C}
\end{align}
\end{lemma}

\begin{proof}
For any $x\neq0$, define
\begin{align}
D(x) &=A(x)-A(-x),\label{Def-D-x}\\
M(x) &=A(x)+\mathscr{S}(x)A(-qx),\label{Def-M-x}\\
N(x) &=\mathscr{S}(x)A(-qx)-A(-x).\label{Def-N-x}
\end{align}
We first derive explicit expressions for these functions. Replacing $q$ by $q^2$ and
setting $c=x$ in \eqref{War-gene}, we find that
\begin{align}
D(x)=2x(-q^2;q^2)_\infty^2\!\:\theta\big({-}a,-b,-abx^2;q^2\big).\label{D-x-exp}
\end{align}
We shall also use the elementary transformation
\begin{align}
\theta\big(q^2x;q^2\big)=-x^{-1}\theta\big(x;q^2\big),\label{Theta-prop}
\end{align}
which follows immediately from the definition of the theta function. Replacing $x$ by
$qx$ in \eqref{D-x-exp}, multiplying both sides by $\mathscr{S}(x)$ and then applying
\eqref{Theta-prop}, we deduce that
\begin{align*}
\mathscr{S}(x)D(qx) &=2qx(-q^2;q^2)_\infty^2\!\:\mathscr{S}(x)\!\:\theta
\big({-}a,-b,-abq^2x^2;q^2\big)
\nonumber\\
 &=2x(-q^2;q^2)_\infty^2\!\:\theta\big({-}a,-b,-abx^2;q^2\big).
\end{align*}
Consequently,
\begin{align}
\mathscr{S}(x)\big(A(qx)-A(-qx)\big)=D(x).\label{Aux-iden-1}
\end{align}

Next, applying \eqref{War-gene} with bases $q$ and $-q$, respectively, yields
\begin{align}
A(x)A(qx)-A(-x)A(-qx) &=2x(-q;q)_\infty^2\!\:\theta\big({-}a,-b,-abx^2;q\big),
\label{Sub-iden-1}\\
A(x)A(-qx)-A(-x)A(qx) &=2x(q;-q)_\infty^2\!\:\theta\big({-}a,-b,-abx^2;-q\big).
\label{Sub-iden-2}
\end{align}
We now rewrite the left-hand sides of \eqref{Sub-iden-1} and \eqref{Sub-iden-2} in
terms of $D$, $M$ and $N$. Indeed,
\begin{align}
 &\mathscr{S}(x)\big(A(x)A(qx)-A(-x)A(-qx)\big)\nonumber\\
 &\qquad\qquad\quad=\mathscr{S}(x)\big(A(qx)-A(-qx)\big)A(x)\nonumber\\
 &\qquad\qquad\qquad\qquad\qquad+\mathscr{S}(x)\big(A(x)-A(-x)\big)A(-qx).
\label{Der-step-1}
\end{align}
It follows from \eqref{Def-D-x}, \eqref{Def-M-x} and \eqref{Aux-iden-1} that
\begin{align}
\mathscr{S}(x)\big(A(x)A(qx)-A(-x)A(-qx)\big)=D(x)M(x).\label{Rel-iden-1}
\end{align}
Similarly,
\begin{align}
\mathscr{S}(x)\big(A(x)A(-qx)-A(-x)A(qx)\big)=D(x)N(x).\label{Rel-iden-2}
\end{align}
Combining \eqref{Sub-iden-1}, \eqref{Sub-iden-2}, \eqref{Rel-iden-1} and
\eqref{Rel-iden-2}, and simplifying the resulting theta products, we obtain that
\begin{align}
M(x) &=(-q;q^2)_\infty^2\!\:\mathscr{S}(x)\!\:\theta\big({-}aq,-bq,-abqx^2;q^2\big),
\label{Aux-iden-2}\\
N(x) &=(q;q^2)_\infty^2\!\:\mathscr{S}(x)\!\:\theta\big(aq,bq,abqx^2;q^2\big).
\label{Aux-iden-3}
\end{align}

The definitions \eqref{Def-D-x}--\eqref{Def-N-x}, together with \eqref{Aux-iden-1},
give
\begin{align}
A(x) &=\dfrac{M(x)+D(x)-N(x)}{2},\label{Rel-iden-3}\\
A(-x) &=\dfrac{M(x)-D(x)-N(x)}{2},\label{Rel-iden-4}\\
\mathscr{S}(x)A(qx) &=\dfrac{M(x)+D(x)+N(x)}{2},\label{Rel-iden-5}\\
\mathscr{S}(x)A(-qx) &=\dfrac{M(x)-D(x)+N(x)}{2}.\label{Rel-iden-6}
\end{align}
Substituting \eqref{Rel-iden-3}--\eqref{Rel-iden-6} into the left-hand side of
\eqref{Fund-iden-1}, we expand the products and collect like terms. Except for the
following six terms, all other terms cancel pairwise:
\begin{align}
 &\mathscr{D}\big(x_1^2,x_2^2,x_3^2\big)-\mathscr{S}\big(x_1^2,x_2^2,x_3^2\big)
\mathscr{D}\big(qx_1^2,qx_2^2,qx_3^2\big)\nonumber\\
 &\qquad\qquad=-\dfrac{1}{2}\sum_{(i,j,k)\in S}D\big(x_i^2\big)\big\{N\big(x_j^2\big)
M\big(x_k^2\big)+M\big(x_j^2\big)N\big(x_k^2\big)\big\},\label{LHS-sum-exp}
\end{align}
where
\begin{align*}
S=\{(1,2,3),\,(2,3,1),\,(3,1,2)\}.
\end{align*}
Finally, substituting \eqref{D-x-exp}, \eqref{Aux-iden-2}, and \eqref{Aux-iden-3} into
\eqref{LHS-sum-exp}, and collecting the common factors, we obtain a linear combination
of theta products. Using the elementary transformations of theta functions, this
expression simplifies to the right-hand side of \eqref{Fund-iden-1}. Hence,
\eqref{Fund-iden-1} follows.

This completes the proof of Lemma~\ref{Fund-Lem-1}.
\end{proof}

\end{document}
